%---------------------------------------------------------------------------%
%->> Chapter 6: 总结与展望
%---------------------------------------------------------------------------%

\chapter{总结与展望}

\section{工作总结}

本文围绕微服务故障根因定位中诊断经验难以有效复用的问题展开研究。针对多源异构观测难以形成统一故障表示、历史诊断经验难以沉淀为可执行知识、在线诊断难以在经验复用与自主探索之间稳定切换等问题，论文从故障表征、知识构建、在线执行和系统实现四个方面展开研究，形成了离线知识构建、在线诊断执行和结果回流更新相衔接的闭环方法框架。本文的主要工作如下。

\begin{enumerate}[label=（\arabic*）]
    \item \textbf{传播感知的多模态故障指纹构建方法。}
    针对指标、日志和调用链路在结构形式、时间尺度与语义层级上的差异，本文首先对多模态观测数据进行统一表示，然后围绕故障传播过程提取传播拓扑、时序关联、日志模式和跨模态耦合等特征，构建可用于相似故障检索的故障指纹。该方法为当前故障与历史故障之间的稳定比较提供了统一表示，也为后续 SOP 的召回与绑定提供了检索入口。

    \item \textbf{面向微服务故障诊断的 SOP 自动化构建方法。}
    针对历史经验分散、结构不统一且难以直接用于智能体执行的问题，本文设计了先验引导的 SOP 执行轨迹生成、决策单元提取与加权、规则组织三个阶段。该方法能够从故障实例与诊断轨迹中抽取判断依据、路径分支和切换条件，并将其组织为可执行的 SOP 结构，使诊断经验由具体案例逐步转化为可复用、可更新的知识表示。

    \item \textbf{SOP 驱动的智能体在线故障诊断方法。}
    在线阶段，系统首先根据当前故障生成故障指纹，并从故障指纹库中召回相似历史案例对应的 SOP。随后，诊断智能体围绕当前焦点推进排查，观察智能体负责采集所需的多模态证据。当现有 SOP 无法覆盖当前分支时，智能体通过逃逸机制扩展新的诊断路径，并在人工复核后将在线诊断记录回流知识库，用于后续故障指纹与 SOP 的增量更新。该方法使历史诊断经验能够在在线诊断中发挥实际作用，同时保留了系统对新故障场景的适应能力。

    \item \textbf{故障诊断智能体平台设计与实验验证。}
    本文设计并实现了一个同时支持离线知识构建与在线故障诊断任务的故障诊断智能体平台，贯通了故障指纹构建、SOP 自动化构建、在线执行和结果回流更新等关键环节。公开数据集上的实验结果表明，本文方法在故障指纹检索、SOP 构建有效性以及端到端根因定位性能方面均取得了较好效果，验证了整体方案的工程可行性和应用价值。
\end{enumerate}

综上，本文围绕历史诊断经验在微服务故障根因定位中的稳定复用问题，完成了从方法设计、系统实现到实验验证的整体研究，构建了面向离线知识沉淀与在线诊断执行的闭环框架，为经验驱动的微服务故障根因定位研究提供了较系统的技术基础。
\section{未来展望}

本文已经完成了方法设计、平台实现与实验验证，但面向真实生产环境的微服务故障诊断仍然存在若干值得进一步研究的问题。后续工作可从以下几个方面展开。

\begin{enumerate}[label=（\arabic*）]
    \item \textbf{面向 SOP 构建过程的效率优化与调用开销控制。}
    本文当前的 SOP 自动化构建依赖于围绕同一故障实例生成多条 SOP 执行轨迹，并基于这些轨迹提取决策单元、计算权重并组织形成有效 SOP。该过程虽然能够较好地恢复诊断步骤并沉淀经验知识，但也带来了较高的构建成本。由于不同执行轨迹通常彼此独立生成，即使它们面向同一故障实例、共享相近的初始状态和部分诊断路径，系统仍需为每条轨迹重复完成一次较完整的根因定位过程，并在过程中触发多轮大语言模型推理与工具调用。随着轨迹数量增加，SOP 构建阶段的模型调用开销会大幅增加，进而限制方法在更大规模知识构建任务中的适用性。后续可围绕状态共享、证据缓存、轨迹前缀复用和中间结果复用等方向展开研究，减少不同轨迹之间的重复推理与重复调用，从而在保证轨迹质量的前提下降低离线知识构建成本，提升 SOP 自动化构建的整体效率。

    \item \textbf{面向更大规模微服务系统的扩展。}
    本文当前的实验验证主要基于公开数据集展开，系统规模与运行复杂度仍相对有限。若进一步面向包含数百甚至上千个服务的生产环境，仍需研究大规模故障指纹构建的开销控制、SOP 图结构的高效检索与匹配，以及复杂服务拓扑下智能体协同执行的可扩展性问题。

    \item \textbf{面向更复杂观测信息的多模态融合。}
    本文当前已经围绕指标、日志和调用链路建立了统一分析框架，但真实运维场景中的观测信息往往还包括监控面板、拓扑可视化图、配置变更记录以及更复杂的长时序模式。后续可结合多模态大模型与时序基础模型，增强系统对复杂观测线索的理解能力，进一步提升故障感知的完整性与诊断推理的精细程度。

    \item \textbf{从根因定位向故障处置延伸。}
    本文的研究重点集中在根因定位阶段，对修复动作和处置策略涉及较少。后续可进一步将修复建议、执行约束与风险控制策略纳入现有知识结构，使系统在完成根因定位之后，还能够围绕候选方案给出处置建议、验证执行条件，并为人工决策提供更直接的支持。

    \item \textbf{提升知识更新过程的自动化与稳健性。}
    本文已经建立了在线诊断结果经人工复核后回流知识库的更新机制，但在知识筛选、置信度评估和错误经验过滤方面仍有进一步优化空间。后续可围绕知识入库标准、案例质量评估和异常更新检测等问题展开研究，以提高知识库持续演化过程中的可靠性。
\end{enumerate}

总体而言，本文构建了一个面向经验驱动微服务根因定位的基础框架。后续工作的重点，在于进一步提升该框架在知识构建效率、系统规模适应能力、复杂观测理解能力以及长期运行稳定性等方面的综合表现，使其更好地适应真实生产环境中的持续演化需求。